\documentclass{stex}

\libusepackage{naproche}
\libusepackage{copyright}
\libinput{preamble}

\begin{document}

\usemodule[libraries/meta]{preliminaries.ftl}


\section{Relations}\label{sec:relations}

\begin{sparagraph}[style=symdoc,for=Eq]
  \begin{signature*}
    Let $x, y$ be terms.
    $\fakeemph{x \Eq y}$ is a relation.
    Let $x$ is \emph{equal to $y$} stand for $x \Eq y$.
    Let $x$ and $y$ \emph{agree} stand for $x \Eq y$.
    Let $x$ \emph{agrees with $y$} stand for $x \Eq y$.
  \end{signature*}
\end{sparagraph}

\begin{sparagraph}[style=symdoc,for=NotEq]
  \begin{definition*}
    Let $x, y$ be terms.
    $\fakeemph{x \NotEq y}$ iff $\NOT x \Eq y$.
    Let $x$ and $y$ are \emph{distinct} stand for $x \NotEq y$.
  \end{definition*}
\end{sparagraph}

\Naproche has a hard-coded induction proof tactic (cf. section
\ref{sec:misc-axioms}) that relies on an ``induction ordering'' $\ILess$.
An important special case of this tactic is the following induction principle:
If, for all $x$, we can show that a predicate $\varphi[y]$ holds for every
$y \ILess x$, then $\varphi[x]$ holds for every $x$.
Typical instances of this principle are the (strong) induction principle from
natural number arithmetics (with $\ILess$ being the $<$-relation) or the
$\in$-induction principle from set theory (with $\ILess$ being the
$\in$-relation).

\begin{sparagraph}[style=symdoc,for=ILess]
  \begin{signature*}
    Let $x, y$ be terms.
    $\fakeemph{x \ILess y}$ is a relation.
    Let $x$ is \emph{inductively less than $y$} stand for $x \ILess y$.
  \end{signature*}

  \begin{axiom*}
    Let $\tau[x_1, \dots, x_n]$ be a term and
    $\psi_1[x_1]$, $\psi_2[x_1,x_2]$, \dots, $\psi_n[x_1, \dots, x_n]$ and $\varphi[x_1, \dots, x_n]$ be formulas.
    Assume that the following holds for all $a_1, \dots, a_n$ with $\psi_i[a_1, \dots, a_i]$ for all $i \Elem \EClass{1, \dots, n}$:
    If $\tau[y_1, \dots, y_n] \ILess \tau[a_1, \dots, a_n]$
    for all $y_1, \dots, y_n$ with $\psi_i[y_1, \dots, y_i]$ for all $i \Elem \EClass{1, \dots, n}$
    then $\varphi[a_1, \dots, a_n]$.
    Then $\varphi[a_1, \dots, a_n]$ for all $a_1, \dots, a_n$ with $\psi_i[a_1, \dots, a_i]$ for all $i \Elem \EClass{1, \dots, n}$.
  \end{axiom*}

  \begin{proposition*}
    Let $A$ be a class, and $\varphi[x]$ be a formula.
    Assume that the following holds for all $a \Elem A$:
    If $\varphi[b]$ for for all $b \Elem A$ with $b \ILess a$ then $\varphi[a]$.
    Then $\varphi[a]$ for all $a \Elem A$.
  \end{proposition*}
\end{sparagraph}


\ifinputref\else\printlicense[CcByNcSa]{Marcel Schütz (2026)}\fi

\end{document}
